\section{Finite Types and Countable Categoricity}
\label{sec:p03-n-a-001-01}

The finite-type criterion connects uniqueness of a countable model to
the number of distinctions expressible about each finite tuple.
The bibliographic title of Ryll-Nardzewski's original article uses
power \(\leq\aleph_0\), including the finite boundary
\cite{RyllNardzewski1959}. Here that boundary is separated from
the infinite-model characterization. The proof develops the
type-realization and extension interfaces of Vaught's formula-algebra
construction (\cref{thm:p03-vaught-finite-formulas-categoricity}),
without identifying this derivation with a different original
author's presentation.

\subsection{The criterion and its hypotheses}
\label{sec:p03-rn-criterion}

\begin{hbconvention}{def:p03-rn-standing}{Complete theories and two meanings of countable categoricity}
In this chapter \(L\) is an at most countable, finite-arity first-order
language with logical identity; constants and function symbols are
allowed. Unless a result explicitly changes the setting, \(T\) is
complete and satisfiable and has an infinite model.
It therefore has only infinite models.
The theory is \(\omega\)-categorical when it has a countably infinite
model and any two countably infinite models are isomorphic.

For the separate finite case, \emph{at-most-countable categoricity}
means existence of a model of size at most \(\aleph_0\) and uniqueness
up to isomorphism among all such models, including finite ones.
These notions agree under the standing infinite-model assumption,
but not without it.

Write \(\bar x=(x_0,\ldots,x_{n-1})\) and let \(S_n(T)\) be the
consistent complete \(n\)-types over the empty set.
Coordinates may repeat and \(n=0\) is included. For
\(\mathcal M\models T\), put
\[
 \operatorname{tp}^{\mathcal M}(\bar a)
       =\{\varphi(\bar x):\mathcal M\models\varphi(\bar a)\}.
\]
The Boolean algebra \(\mathbb B_n(T)\) is the quotient of formulas
in these slots modulo equivalence in every \(T\)-model, as in
\cref{def:p03-vaught-formula-algebra}.
An isolator is consistent with \(T\) and implies every formula
of its complete type; the consistency condition is essential.
The constructions use the compactness and Choice standing of
\cref{thm:fol-compactness}.
\end{hbconvention}

\begin{hblemma}{thm:p03-rn-type-realization}{proved=annex:p3.rn.type-realization}{Every complete finite-tuple type has a countable realization}
For each finite \(n\) and \(p\in S_n(T)\), a countably infinite
\(T\)-model realizes \(p\).
Named coordinates may coincide; no distinctness is imposed unless
the type requires it. The unique zero-type is realized by the empty
tuple. The result includes countable languages with functions.
\end{hblemma}

\begin{hblemma}{thm:p03-rn-isolated-extension}{proved=annex:p3.rn.isolated-extension}{Transfer of an isolated joint type over a fixed tuple}
Suppose \(p\in S_{n+1}(T)\) is isolated by \(\theta(\bar x,y)\),
and \(\bar a\) in a \(T\)-model realizes its first-\(n\)-slot
restriction. Some \(b\) in that same model realizes \(p\)
together with \(\bar a\).
If all empty-parameter finite-tuple types are isolated, every
model realizes all complete finite-tuple types over each finite
parameter set in that model.
\end{hblemma}

\begin{hbtheorem}{thm:p03-rn-characterization}{proved=annex:p3.rn.characterization}{The complete finite-type characterization}
Under \cref{def:p03-rn-standing}, the following are equivalent:
\begin{enumerate}
\item \(T\) is \(\omega\)-categorical.
\item \(S_n(T)\) is finite for every \(n<\omega\).
\item Every \(p\in S_n(T)\), for every finite \(n\), is isolated.
\item \(\mathbb B_n(T)\) is finite for every finite \(n\).
\item Every model of \(T\) is atomic.
\item Every model of \(T\) is \(\aleph_0\)-saturated:
      it realizes types over finite parameter sets.
\item Every countably infinite model of \(T\) is
      \(\aleph_0\)-saturated.
\item \(T\) has a countably infinite model that is both prime
      and \(\aleph_0\)-saturated.
\end{enumerate}
The finite number \(|S_n(T)|\) may depend on \(n\);
no single bound for all tuple lengths is required.
There is no finite-signature or function-free restriction
on this general theorem.
\end{hbtheorem}

The uniqueness-to-isolation obstruction is a type realized
in one countable model and omitted in another. In the reverse
direction, isolation permits an extension over each already
matched finite tuple. Alternating extensions in two enumerated
models gives an onto elementary map. Neither mere density
of isolated types nor a single successful tuple extension
would suffice for this criterion.

\begin{hbcorollary}{thm:p03-rn-orbits}{proved=annex:p3.rn.orbits}{Types are the finite-tuple automorphism orbits}
If \(T\) satisfies the equivalent conditions and
\(\mathcal M\models T\) is countably infinite, two finite
tuples of the same length are in the same
\(\operatorname{Aut}(\mathcal M)\)-orbit exactly when
they have the same complete type.
Every type in \(S_n(T)\) is realized in \(\mathcal M\),
so there are exactly \(|S_n(T)|\) orbits on its \(n\)-tuples.
This includes the one orbit on zero-tuples.
The construction extends the prescribed tuple map.
\end{hbcorollary}

\subsection{Finite, incomplete and uncountable-language boundaries}
\label{sec:p03-rn-boundaries}

\begin{hbtheorem}{thm:p03-rn-finite-boundary}{proved=annex:p3.rn.finite-boundary}{The finite complete-theory alternative}
Let \(T\) be complete and satisfiable in an at most countable
language and have a nonempty finite model of size \(m\).
Every \(T\)-model is isomorphic to it.
Thus \(T\) is at-most-countable categorical, but not
\(\omega\)-categorical under the existence convention above.
For every finite \(n\),
\[
 |S_n(T)|\leq m^n,\qquad
 |\mathbb B_n(T)|\leq2^{m^n}.
\]
Every complete type is isolated and realized in this model,
which is atomic, prime and saturated over every parameter
set in its finite domain. A finite-type criterion alone
does not imply existence of an infinite model.
\end{hbtheorem}

\begin{hbtheorem}{thm:p03-rn-incomplete-boundary}{proved=annex:p3.rn.incomplete-boundary}{Infinite-model uniqueness does not replace completeness}
In the language with only logical identity, the empty theory
has isomorphic countably infinite models but infinitely many
complete one-types: each positive finite exact-cardinality
sentence belongs to a different type.
It also has models of every positive finite size.
Thus uniqueness at the countably infinite cardinal, without
completeness or an assumption that all models are infinite,
does not imply the finite-type criterion.
\end{hbtheorem}

\begin{hbtheorem}{thm:p03-rn-uncountable-language}{proved=annex:p3.rn.uncountable-language}{A genuine countable-language boundary}
There is a complete theory in an uncountable language with
exactly one at most countable model, which is infinite,
but infinitely many complete one-types.
Take the complete theory of \(\omega\) with constants \(c_i\)
for all \(i<\omega\) and a unary function \(f_g\) for every
\(g:\omega\to\omega\), interpreted by \(g\).
Every model containing an element outside all the named
constants has at least \(2^{\aleph_0}\) elements.
The language's countability is an actual hypothesis in
the omission-based converse.
\end{hbtheorem}

\begin{hbconvention}{def:p03-rn-parameter-boundary}{Finite parameters are not countable parameters}
The criterion's \(\aleph_0\)-saturation concerns finite
parameter sets, not arbitrary countable ones.
In any countably infinite model, the consistent type asking
for an element different from every member of its domain
is omitted over that whole domain.
The type count is over the empty set of a fixed language;
infinitely many new names can destroy it.
The rationally named dense order and its irrational-cut
types in \cref{thm:p03-vaught-dlo-example} give the explicit
example.
\end{hbconvention}

\subsection{Dense linear order: the full type calculation}
\label{sec:p03-rn-dlo}

\begin{hbdefinition}{def:p03-rn-dlo-data}{Dense order without endpoints and finite order patterns}
The language here has just \(<\) and logical identity:
no named points, extra predicates or function symbols.
Let \(T_{\rm dlo}\) consist of the strict linear-order axioms,
density
\[
 \forall x\forall y\bigl(x<y\to
                     \exists z(x<z\land z<y)\bigr),
\]
and the absence of both endpoints
\[
 \forall x\exists y(y<x),\qquad
 \forall x\exists y(x<y).
\]
An order pattern is an ordered partition \((B_1,\ldots,B_k)\)
of \(\{0,\ldots,n-1\}\) into nonempty blocks.
Its formula \(\delta_{(B_1,\ldots,B_k)}(\bar x)\)
requires equal coordinates within each block and
\(x_i<x_j\) whenever \(i\in B_r\), \(j\in B_s\), \(r<s\).
There is one empty pattern for \(n=0\), with formula \(\top\).
The number \({n\brace k}\) counts partitions into exactly
\(k\) nonempty unordered blocks; \({0\brace0}=1\).
\end{hbdefinition}

\begin{hblemma}{thm:p03-rn-dlo-extension}{proved=annex:p3.rn.dlo-extension}{Every finite order map extends in both directions}
An order-preserving bijection between finite subsets of two
\(T_{\rm dlo}\)-models extends to any further specified
element of either domain.
An old element uses its existing mate.
A new one is matched below the first old point, above the last,
or between the appropriate adjacent old points.
With an empty initial subset any target point suffices.
Each open interval and end ray is infinite, so a fresh target
can be selected when needed.
No countability hypothesis is used in this finite extension.
\end{hblemma}

\begin{hbtheorem}{thm:p03-rn-dlo-pattern-truth}{proved=annex:p3.rn.dlo-pattern-truth}{Order patterns decide every formula}
Tuples in any two \(T_{\rm dlo}\)-models with the same order
pattern satisfy the same first-order formulas in their slots.
Thus \(T_{\rm dlo}\) is complete, and each formula is
equivalent modulo it to the finite disjunction of the order
patterns making it true.
This proves quantifier elimination without assuming it.
Every pattern is realized in every model.
\end{hbtheorem}

\begin{hbcorollary}{thm:p03-rn-dlo-type-count}{proved=annex:p3.rn.dlo-type-count}{Ordered partitions are exactly the complete types}
The pattern formulas isolate exactly the complete types.
Therefore
\[
 |S_n(T_{\rm dlo})|=\sum_{k=0}^n k!\,{n\brace k},
 \qquad
 |\mathbb B_n(T_{\rm dlo})|=2^{\,|S_n(T_{\rm dlo})|}.
\]
The type counts for \(n=0,1,2,3\) are \(1,1,3,13\).
The two-coordinate types are equality and the two strict
orderings; the thirteen three-coordinate types are one
all-equal pattern, six two-block patterns and six three-block
patterns.
\end{hbcorollary}

\begin{hbtheorem}{thm:p03-rn-dlo-uniqueness}{proved=annex:p3.rn.dlo-uniqueness}{The rational model and prescribed-map back-and-forth}
The ordered rationals form a countably infinite
\(T_{\rm dlo}\)-model. Every countably infinite dense linear
order without endpoints is isomorphic to it by alternating
finite order extensions.
The construction extends any prescribed finite order map
between two such models.
Consequently the rational model is homogeneous, its tuple
orbits are precisely the ordered partitions, and
\(T_{\rm dlo}\) is \(\omega\)-categorical.
Density without both endpoint clauses would not justify
the end-ray extension steps.
\end{hbtheorem}

\subsection{The corrected finite-relational application}
\label{sec:p03-rn-finite-relational}

The general theorem allows functions and infinite countable
signatures. A useful \emph{application} has stronger hypotheses:
Hodges's author correction to the finite-signature application
at p.~352 adds the absence of function symbols
\cite[p.~3, correction to p.~352]{Hodges1993Corrigenda2011}.
The following argument makes that restriction's role explicit.

\begin{hbtheorem}{thm:p03-rn-relational-application}{proved=annex:p3.rn.relational-application}{Finite relational homogeneity implies countable categoricity}
Let \(\mathcal M\) be countably infinite in a finite,
finite-arity, purely relational language without functions
or constants. Suppose every isomorphism between finite
induced substructures extends to an automorphism.
Then \(\operatorname{Th}(\mathcal M)\) has quantifier
elimination and is \(\omega\)-categorical.
For each fixed \(n\), the finitely many atomic relation
and equality patterns give a finite list of complete types.
Finiteness of this list and extendibility of the finite maps
are separate steps.
\end{hbtheorem}

\begin{hbtheorem}{thm:p03-rn-function-boundary}{proved=annex:p3.rn.function-boundary}{Finite function signature does not imply a finite pattern list}
Let \(\mathcal Z=(\mathbb Z,s,t)\), with
\(s(a)=a+1\) and \(t(a)=a-1\).
Every nonempty finitely generated substructure is the whole
structure; an isomorphism between such substructures is
already an automorphism.
Nevertheless its complete theory is not
\(\omega\)-categorical: the pairs \((0,k)\), \(k\geq0\),
have distinct two-types, distinguished by \(y=s^k(x)\).
There is no nonempty finite closed substructure.
This example does not satisfy the finite induced-substructure
hypothesis; it shows why finite function signatures cannot
silently be treated as finite relational pattern calculations.
\end{hbtheorem}

The source identities have distinct extents. The selected
Engeler 1959 entry is a Notices announcement at p.~161,
not a longer proof article \cite{Engeler1959}.
The selected Svenonius categoricity article is Theoria 25,
pp.~82--94 \cite{Svenonius1959}, distinct from his same-year
permutation article at pp.~173--178.
The Hodges correction concerns the bounded application,
not a new function-free hypothesis on the general criterion.
